% sec-09-conclusions.tex - Section 9, Conclusions.  FINAL stage.
% One page.  No figures and no tables: everything here has already been shown,
% and the conclusion's job is to state what follows from it.  Senior-author
% pass: the two section cross-references are made symbolic so they survive a
% renumber, and a closing paragraph states the incompatibility claim without
% criticizing the practitioners of the prior system.
\section{Conclusions}
\label{sec:conclusions}

Between January 2025 and July 2026 the executive branch built the demand side of
an AI-assisted cancer research program and left the supply side alone. Executive
Order 14179 removed barriers to Federal AI adoption and Executive Order 14212
created the commission that Executive Order 14355 then pointed at cancer
\cite{eo14179, eo14212, eo14355}. The Childhood Cancer Data Initiative budget
doubled to \$100 million on the day that order was signed \cite{ccdidoubling_nih}.
The Genesis Mission and its NIH arm committed more than \$1.2 billion
\cite{eo14363, biogenesis_nih}. The FDA fielded a generative tool for tasks
including clinical protocol review and followed it with an agency-wide agentic
platform \cite{crs_ai_healthcare, hhs_ai_strategy}. CMS built an
interoperability framework and a technology office to hold it
\cite{cms_healthtech, cms_office_healthtech}. On June 22, 2026 HHS launched a
department-wide effort to restore American leadership in clinical trials, with
Operation TrialBlazer as its oncology arm \cite{trialblazer, hhstrialblazer2026}.
Compute, data, agency review capacity, and faster site activation are all being
funded. The document a sponsor must file is not, and that is the gap this paper
closes on the sponsor side, with artifacts rather than with a proposal.

An IND of twelve modules under 21 CFR \S 312 and \S 812, carrying 594,657
characters, 22 figures, 84 section-scoped tables and 52 codified
cross-references, was assembled in four days by one author \cite{phase1ind}. A
Phase 1 protocol at n = 18 behind a Phase 0 gate of at least 1000 simulations
across two independent frameworks, and a Phase 2 protocol at n = 220 across
eight high-volume centers targeting a hazard ratio of 0.60, were deposited two
days apart and total 369,299 characters and 42 figures
\cite{onpremwhippl, phase2proto}. Five bill versions and four companion
documents were deposited across nineteen days, ending in a financial data
amendment that attaches a cost ledger to every verification run
\cite{hr9510billv5, clintrust, congressvote}. Fourteen funding artifacts
followed, carrying a \$1,606,000 small-business ask against a \$3,500,000
five-year direct program, \$5,900,000 of private capital above a firewall at
3.67 to one, and a virtual trial projected at \$36,330 against an industry
benchmark above \$120,000
\cite{capplan2026, tenapps2026, rfarm27001v2}. Every one of those artifacts was
reviewed during development by two model manufacturers independent of the one
that produced it, and cleared by a human principal investigator who retained
final authority \cite{aiprgliomatc, rfarm27001v2}.

The two systems are primarily incompatible, and the incompatibility is not a
matter of preference. A review regime whose best case is seven to eight weeks
and whose typical round is several months cannot absorb an author who deposits
more than thirty artifacts in a year, and a review that arrives after release
cannot correct the object that was released. A production regime whose unit is
the team-month cannot be audited by a commit history it does not keep. A figure
format that stores pixels cannot be read back as the input to the next paper.
Each of those three is a property of the prior system's architecture rather than
of its practitioners, and none is fixable by working harder inside it. The new
system's advantages are correspondingly persistent: comprehensive provenance on
every artifact, a 1 to 4 day project time scale, AI peer review during
development rather than after it, and machine-readable figures that are re-read
rather than re-drawn.

What follows is a funding decision, not another document. The author is seeking
near-term funding for the PDAC hybrid clinical trial whose IND and protocols are
described in \S\ref{sec:ind} and \S\ref{sec:protocol}, and every artifact named in this paper exists, is
deposited under a digital object identifier, and can be read before any money
moves. Nothing here is an approved protocol, an active IND, or an enacted
statute, and this paper does not claim otherwise: what it claims is that the
documents a sponsor must produce can now be produced correctly, comprehensively,
and fast enough that the production of them is no longer the rate-limiting step
in bringing a pancreatic cancer trial to a patient.

One sentence is worth leaving with a funder. The prior system's rate-limiting
step was the document; in the new system it is the institution, and the
institution can be bought with money in a way the document could not be bought
with effort.

The prior system is not being criticized for lacking a capability it could have
had. It is being described as an architecture whose throughput was matched to a
production rate that no longer holds on one side of the transaction, and the
right response to that is not to defend either side but to say plainly which
parts are now incompatible and to build accordingly. This paper has done that
across five main sections, twenty-five figures and twenty-five tables, and every
one of the artifacts it reports on is deposited and readable today.
